\chapter{Resultados Preliminares}

O retificador PFC trifásico de ponte completa é simulado no PSIM com base nos parâmetros
obtidos por \cite{tech_simulator_2020}, os quais são listados abaixo:

\begin{itemize}
    \item Tensão de linha: \quad \(V_\mathrm{ff} = 415\ \mathrm{V_{rms}}\)
    \item Frequência da rede elétrica: \quad \(f = 60\ \mathrm{Hz}\)
    \item Indutor de entrada: \quad \( L = 1\ \mathrm{mH}\)
    \item Capacitor de saída: \quad \( L = 470\ \mathrm{\mu F}\)
    \item Resistor de saída: \quad \( R = 6,4\ \mathrm{\Omega}\)
\end{itemize}

O diagrama de simulação do retificador PFC trifásico é visto na figura
\ref{fig:diagrama_retificador_pfc}. As tensões e correntes das fases a,b,c obtidas na simulação
são exibidas na figura \ref{fig:diagrama_tensoes_correntes}.

\begin{figure}[h]
    \centering
    \caption{Diagrama de simulação do retificador PFC trifásico no PSIM.}
    \includegraphics[width=0.8\textwidth]{figuras/Retificador3ph_PSIM.png}
    \legend{Fonte: Autoral}
    \label{fig:diagrama_retificador_pfc}
\end{figure}

\newpage

\begin{figure}[h]
    \centering
    \caption{Gráfico das correntes e tensões da simulação do retificador PFC trifásico.}
    \includegraphics[width=0.8\textwidth]{figuras/three_phase_plot.png}
    \legend{Fonte: Autoral}
    \label{fig:diagrama_tensoes_correntes}
\end{figure}
